\textbf{Capacidad seleccionada y escalado.}
La línea base utiliza ocho unidades desplegables: tres núcleos y cinco procesos especializados. Cada unidad mantiene dos tareas mínimas en dos zonas; sincronización utiliza 2 vCPU/8 GiB por tarea y las otras unidades 1 vCPU/4 GiB. El escenario normal suma $3(2)+5(2)=16$ tareas, 18 vCPU y 72 GiB; peak utiliza 28 tareas, 32 vCPU y 128 GiB; crecimiento, 46 tareas, 52 vCPU y 208 GiB. El techo estable es de ocho tareas por unidad: 64 tareas, 72 vCPU y 288 GiB. El Anexo A4.6 conserva cantidades por unidad, fórmulas, reservas y diferencias de ambientes \parencite{c8-fargate-recursos}.

Los escenarios corresponden a 25, 100 y 300 HTTP/s y a la mezcla concurrente de SD5, 5.4.1. Como hipótesis de diseño, el núcleo más cargado recibe 50\,\% de solicitudes, demanda 0,010 segundos de CPU por solicitud y reserva 60\,\% de CPU para negocio: en crecimiento requiere $300(0,50)(0,010)=1,50$ vCPU. Tres tareas supervivientes aportan $3(0,60)=1,80$ vCPU. La diferencia es un margen calculado; no constituye rendimiento medido. A4.6 incorpora agentes, trabajos pesados, integración, notificaciones y sincronización para evitar asignar el mismo margen varias veces.

Los núcleos escalan con CPU de aplicación objetivo de 60\,\%, memoria total de 70\,\% y latencia P95 por tipo de operación. Los workers agregan tareas ante edad de cola mayor de treinta segundos o más de cien pendientes por tarea durante dos minutos. Se amplía con enfriamiento de sesenta segundos y se reduce después de quince minutos; no se reduce durante emergencia, reconexión, despliegue o pérdida de zona. Alcanzar el máximo con cola o latencia creciente exige ampliar la línea base antes de habilitar más carga \parencite{c8-ecs-autoescalado}.

Se seleccionan dos almacenes RDS PostgreSQL 18 Multi-AZ de 2 vCPU/8 GiB, con 200 y 50 GiB gp3 en base. Crecimiento contempla 8 vCPU/32 GiB y al menos 400 GiB por almacén; DR se amplía antes de su primario. IOPS y caudal de volumen no prueban el rendimiento de la aplicación. Desarrollo y QA usan una tarea por unidad; Preproducción reproduce recursos, topología y políticas del escenario ensayado; DR conserva dieciséis tareas mínimas y adopta la capacidad productiva antes de promoverse \parencites{c8-rds-clases,c8-rds-almacenamiento}.

Cada promoción exige carga y estrés a $1,5$ veces el peak habilitado: 150 HTTP/s para peak inicial de 100, o 450 si se habilitan 300, junto con eventos, avisos y reconexión. Se prueban pérdida de tarea/zona, integridad, colas, latencias y recuperación. SRE provisiona el escenario y Calidad verifica sus resultados; las diferencias de Preproducción se justifican en EQ-4. La matriz y las fórmulas de A4.6 sustentan la aceptación futura, sin atribuir ensayos ejecutados ni incluir montos de la Oferta Económica.
